\clearpage
\begin{proofbox}
\textbf{\textcolor{CProof}{方法一：任取两个数比较（最推荐，初中不等式知识就够）}}
\begin{center}\begin{tikzpicture}[font=\small,>=stealth,line join=round]
\draw[gray!40,rounded corners=3pt] (0,0) rectangle (14.2,14.0);
\node at (7.1,13.4) {四种组合：看同一对输入的先后顺序怎样变化};
\begin{scope}[shift={(0.7,7.5)}]\node at (3.0,5.0) {里增、外增：整体增};
\draw[->,gray] (.4,4.0)--(5.7,4.0);\node[left] at (.4,4.0) {$x$};\fill[blue] (1.35,4.0) circle(2pt);\draw[blue,dashed,->,shorten <=3pt,shorten >=3pt] (1.35,4.0)--(1.35,2.5);\fill[orange!85!black] (4.65,4.0) circle(2pt);\draw[orange!85!black,dashed,->,shorten <=3pt,shorten >=3pt] (4.65,4.0)--(4.65,2.5);\draw[->,gray] (.4,2.5)--(5.7,2.5);\node[left] at (.4,2.5) {$t$};\fill[blue] (1.35,2.5) circle(2pt);\draw[blue,dashed,->,shorten <=3pt,shorten >=3pt] (1.35,2.5)--(1.35,1.0);\fill[orange!85!black] (4.65,2.5) circle(2pt);\draw[orange!85!black,dashed,->,shorten <=3pt,shorten >=3pt] (4.65,2.5)--(4.65,1.0);\draw[->,gray] (.4,1.0)--(5.7,1.0);\node[left] at (.4,1.0) {$y$};\fill[blue] (1.35,1.0) circle(2pt);\fill[orange!85!black] (4.65,1.0) circle(2pt);\node at (3,.25) {$g(x)= x,\ f(t)= t$};\end{scope}
\begin{scope}[shift={(7.5,7.5)}]\node at (3.0,5.0) {里减、外减：整体增};
\draw[->,gray] (.4,4.0)--(5.7,4.0);\node[left] at (.4,4.0) {$x$};\fill[blue] (1.35,4.0) circle(2pt);\draw[blue,dashed,->,shorten <=3pt,shorten >=3pt] (1.35,4.0)--(4.65,2.5);\fill[orange!85!black] (4.65,4.0) circle(2pt);\draw[orange!85!black,dashed,->,shorten <=3pt,shorten >=3pt] (4.65,4.0)--(1.35,2.5);\draw[->,gray] (.4,2.5)--(5.7,2.5);\node[left] at (.4,2.5) {$t$};\fill[blue] (4.65,2.5) circle(2pt);\draw[blue,dashed,->,shorten <=3pt,shorten >=3pt] (4.65,2.5)--(1.35,1.0);\fill[orange!85!black] (1.35,2.5) circle(2pt);\draw[orange!85!black,dashed,->,shorten <=3pt,shorten >=3pt] (1.35,2.5)--(4.65,1.0);\draw[->,gray] (.4,1.0)--(5.7,1.0);\node[left] at (.4,1.0) {$y$};\fill[blue] (1.35,1.0) circle(2pt);\fill[orange!85!black] (4.65,1.0) circle(2pt);\node at (3,.25) {$g(x)=-x,\ f(t)=-t$};\end{scope}
\begin{scope}[shift={(0.7,1.5)}]\node at (3.0,5.0) {里增、外减：整体减};
\draw[->,gray] (.4,4.0)--(5.7,4.0);\node[left] at (.4,4.0) {$x$};\fill[blue] (1.35,4.0) circle(2pt);\draw[blue,dashed,->,shorten <=3pt,shorten >=3pt] (1.35,4.0)--(1.35,2.5);\fill[orange!85!black] (4.65,4.0) circle(2pt);\draw[orange!85!black,dashed,->,shorten <=3pt,shorten >=3pt] (4.65,4.0)--(4.65,2.5);\draw[->,gray] (.4,2.5)--(5.7,2.5);\node[left] at (.4,2.5) {$t$};\fill[blue] (1.35,2.5) circle(2pt);\draw[blue,dashed,->,shorten <=3pt,shorten >=3pt] (1.35,2.5)--(4.65,1.0);\fill[orange!85!black] (4.65,2.5) circle(2pt);\draw[orange!85!black,dashed,->,shorten <=3pt,shorten >=3pt] (4.65,2.5)--(1.35,1.0);\draw[->,gray] (.4,1.0)--(5.7,1.0);\node[left] at (.4,1.0) {$y$};\fill[blue] (4.65,1.0) circle(2pt);\fill[orange!85!black] (1.35,1.0) circle(2pt);\node at (3,.25) {$g(x)= x,\ f(t)=-t$};\end{scope}
\begin{scope}[shift={(7.5,1.5)}]\node at (3.0,5.0) {里减、外增：整体减};
\draw[->,gray] (.4,4.0)--(5.7,4.0);\node[left] at (.4,4.0) {$x$};\fill[blue] (1.35,4.0) circle(2pt);\draw[blue,dashed,->,shorten <=3pt,shorten >=3pt] (1.35,4.0)--(4.65,2.5);\fill[orange!85!black] (4.65,4.0) circle(2pt);\draw[orange!85!black,dashed,->,shorten <=3pt,shorten >=3pt] (4.65,4.0)--(1.35,2.5);\draw[->,gray] (.4,2.5)--(5.7,2.5);\node[left] at (.4,2.5) {$t$};\fill[blue] (4.65,2.5) circle(2pt);\draw[blue,dashed,->,shorten <=3pt,shorten >=3pt] (4.65,2.5)--(4.65,1.0);\fill[orange!85!black] (1.35,2.5) circle(2pt);\draw[orange!85!black,dashed,->,shorten <=3pt,shorten >=3pt] (1.35,2.5)--(1.35,1.0);\draw[->,gray] (.4,1.0)--(5.7,1.0);\node[left] at (.4,1.0) {$y$};\fill[blue] (4.65,1.0) circle(2pt);\fill[orange!85!black] (1.35,1.0) circle(2pt);\node at (3,.25) {$g(x)=-x,\ f(t)= t$};\end{scope}

\node[text width=12.8cm,align=center] at (7.1,.65) {各数轴相同比例，坐标为 $-1,1$；蓝点追踪 $x_1=-1$，橙点追踪 $x_2=1$。};
\end{tikzpicture}\end{center}



\textbf{为什么选这个方法？} 判断一个函数是不是递增，按定义就是：
在区间内随便取两个数 $x_1<x_2$，看后者的函数值是否总比前者大。
这个方法\textbf{不用求导、没有额外假设}，最能解释口诀从哪里来。

\textbf{第一步：}设 $H(x)=f(g(x))$，任取 $x_1,x_2\in I$，且 $x_1<x_2$。
记 $t_1=g(x_1)$，$t_2=g(x_2)$。
因为 $g$ 在 $I$ 上严格单调，所以 $t_1,t_2$ 的大小一定能确定；
又因为这两个数都落在 $f$ 保持同一增减方向的范围内，
可以继续比较 $f(t_1)$ 和 $f(t_2)$。

\end{proofbox}
\clearpage
\begin{proofbox}
\textbf{第二步：逐种比较。}
\begin{enumerate}[leftmargin=*, itemsep=0.5em]
\item \textbf{里增、外增。} 由 $x_1<x_2$ 得 $t_1<t_2$；
外层递增，所以 $f(t_1)<f(t_2)$，即 $H(x_1)<H(x_2)$。整体递增。
\[
 x_1<x_2\ \Longrightarrow\ g(x_1)<g(x_2)
 \ \Longrightarrow\ f(g(x_1))<f(g(x_2)).
\]
\item \textbf{里减、外减。} 内层先把不等号反过来，得到 $t_1>t_2$；
外层也递减，因此输入更大的 $t_1$，输出反而更小：$f(t_1)<f(t_2)$。
整体仍递增。
\[
 x_1<x_2\ \Longrightarrow\ g(x_1)>g(x_2)
 \ \Longrightarrow\ f(g(x_1))<f(g(x_2)).
\]
\item \textbf{里增、外减。} 先得到 $t_1<t_2$，外层递减会把不等号反过来：
$f(t_1)>f(t_2)$。整体递减。
\[
 x_1<x_2\ \Longrightarrow\ g(x_1)<g(x_2)
 \ \Longrightarrow\ f(g(x_1))>f(g(x_2)).
\]
\item \textbf{里减、外增。} 先得到 $t_1>t_2$，外层递增不会改变这两个
$t$ 对应的大小顺序：$f(t_1)>f(t_2)$。整体递减。
\[
 x_1<x_2\ \Longrightarrow\ g(x_1)>g(x_2)
 \ \Longrightarrow\ f(g(x_1))>f(g(x_2)).
\]
\end{enumerate}
\textbf{第三步：}由于 $x_1,x_2$ 是在区间内\textbf{任意}取的，
结果就对整段区间都成立。这才是证明，不是只试了几个数。

\medskip
\end{proofbox}
\clearpage
\begin{proofbox}
\textbf{\textcolor{CProof}{方法二：看两次“变化量”的正负（只用初中代数）}}
\begin{center}\begin{tikzpicture}[font=\small,>=stealth,line join=round]
\draw[gray!40,rounded corners=3pt] (0,0) rectangle (14.2,12.8);
\node at (7.1,12.200000000000001) {两个差商相乘，就是整体的差商};
\begin{scope}[shift={(2.2,8.1)},x=.95cm,y=.95cm]\draw[->,gray] (-.5,0)--(3.7,0);\draw[->,gray] (0,-0.3)--(0,3.8);\draw[blue,thick,domain=-.3:1.4] plot(\x,{2*\x+1});\draw[orange!85!black,dotted] (0,1)--(1,1)--(1,3);\fill[orange!85!black] (0,1) circle(2pt);\fill[orange!85!black] (1,3) circle(2pt);\end{scope}\node[align=left,anchor=west] at (8.4,8.799999999999999) {$g(x)=2x+1$\\[9pt]内层差商 $\dfrac{3-1}{1-0}=2$};
\begin{scope}[shift={(2.2,6.0)},x=.95cm,y=.95cm]\draw[->,gray] (-.5,0)--(3.7,0);\draw[->,gray] (0,-1.9)--(0,0.4);\draw[blue,thick,domain=0:3.5] plot(\x,{-.5*\x});\draw[orange!85!black,dotted] (1,-0.5)--(3,-0.5)--(3,-1.5);\fill[orange!85!black] (1,-0.5) circle(2pt);\fill[orange!85!black] (3,-1.5) circle(2pt);\end{scope}\node[align=left,anchor=west] at (8.4,6.7) {$f(t)=-t/2$\\[9pt]外层差商 $\dfrac{-3/2+1/2}{3-1}=-1/2$};
\begin{scope}[shift={(2.2,3.3)},x=.95cm,y=.95cm]\draw[->,gray] (-.5,0)--(3.7,0);\draw[->,gray] (0,-1.9)--(0,0.4);\draw[blue,thick,domain=-.3:1.5] plot(\x,{-\x-.5});\draw[orange!85!black,dotted] (0,-0.5)--(1,-0.5)--(1,-1.5);\fill[orange!85!black] (0,-0.5) circle(2pt);\fill[orange!85!black] (1,-1.5) circle(2pt);\end{scope}\node[align=left,anchor=west] at (8.4,4.0) {$H(x)=-x-1/2$\\[9pt]整体差商 $-1=2\times(-1/2)$};

\node[text width=12.8cm,align=center] at (7.1,.65) {三图横纵等比例；两点差商依次为 $2,-1/2,-1$。};
\end{tikzpicture}\end{center}



这种写法把口诀变成一条能算出来的式子，不需要导数。
仍任取 $x_1<x_2$，记 $t_i=g(x_i)$，$H(x)=f(g(x))$。
因为 $g$ 严格单调，$t_2-t_1\ne0$。把下式右边乘开：
\[
\boxed{\frac{H(x_2)-H(x_1)}{x_2-x_1}
 =\underbrace{\frac{f(t_2)-f(t_1)}{t_2-t_1}}_{\text{看外层}}
 \cdot
 \underbrace{\frac{t_2-t_1}{x_2-x_1}}_{\text{看内层}}.}
\]
右式约去非零的 $t_2-t_1$ 就得到左式；这里是两点的差商，不是导数。

\textbf{再看正负：}
\begin{itemize}[leftmargin=*, itemsep=0.3em]
\item 内层递增，则右边第二个分式为正；内层递减则为负。
\item 外层递增，则右边第一个分式为正；外层递减则为负。
\item 正数乘正数或负数乘负数，都是正数；正负相乘是负数。
\end{itemize}
而 $x_2-x_1>0$，所以左边分式的正负就是 $H(x_2)-H(x_1)$ 的正负。
于是两层方向相同则整体递增，不同则整体递减。

\medskip
\end{proofbox}
\clearpage
\begin{proofbox}
\textbf{\textcolor{CProof}{方法三：学过导数以后，一行判断（高中方法，可选）}}
\begin{center}\begin{tikzpicture}[font=\small,>=stealth,line join=round]
\draw[gray!40,rounded corners=3pt] (0,0) rectangle (14.2,11.3);
\node at (7.1,10.700000000000001) {例2：外层递减，内层最低点变成整体最高点};

\begin{scope}[shift={(2,7.3)},x=1.2cm,y=.55cm]
\draw[->,gray] (-1.5,0)--(3.6,0) node[right]{$x$};\draw[->,gray] (0,-1.6)--(0,4.5) node[above]{$t$};
\draw[blue,very thick,domain=-1:3,samples=100] plot(\x,{(\x-1)^2-1});\fill[orange!85!black] (1,-1) circle(2pt);\draw[gray,dotted] (1,-1)--(1,3.7);
\end{scope}
\begin{scope}[shift={(2,2.1)},x=1.2cm,y=1.0cm]
\draw[->,gray] (-1.5,0)--(3.6,0) node[right]{$x$};\draw[->,gray] (0,0)--(0,3.7) node[above]{$H$};
\draw[violet,very thick,domain=-1.3:3.3,samples=150] plot(\x,{pow(3,1-(\x-1)^2)});\fill[orange!85!black] (1,3) circle(2pt);\draw[gray,dotted] (1,0)--(1,3);\node[below=5pt] at (1,0){$1$};
\end{scope}
\node[align=left,anchor=west] at (8.3,8.1){$g(x)=(x-1)^2-1$\\[12pt]$x=1$：内层最小 $-1$\\[12pt]$g'(x)=2(x-1)$};
\node[align=left,anchor=west] at (8.3,3.9){$H(x)=3^{-g(x)}$\\[12pt]外层 $3^{-t}$ 递减\\[12pt]内层的谷，对应整体的峰\\[12pt]$H(1)=3$ 为最大值};

\node[text width=12.8cm,align=center] at (7.1,.65) {两图共享相同的 $x$ 刻度；纵轴分别为内层 $t$ 与整体函数值 $H$。};
\end{tikzpicture}\end{center}



\textbf{使用前提：}内层 $g$ 和外层 $f$ 在有关位置都能求导，且区间内相关导数
始终为正或始终为负。如果有导数等于 $0$ 的点，不能单凭“全程大于或小于 $0$”这条简便判断，
需要结合具体情况；\textbf{不能把“导数某点为零”误认为函数一定掉头。}

对于 $H(x)=f(g(x))$，求导时先对外层求导，里面保留 $g(x)$，
再乘里面的导数：
\[
\boxed{H'(x)=f'(g(x))\,g'(x).}
\]
当 $g'(x)>0$，说明里面随 $x$ 增大而增大；$g'(x)<0$ 则相反。
同样，$f'(g(x))>0$ 说明外层在这些输入处递增，$<0$ 说明递减。
因此
\[
\begin{array}{c|c|c}
 g'(x)&f'(g(x))&H'(x)\\ \hline
 +&+&+\\ -&-&+\\ +&-&-\\ -&+&-
\end{array}
\]
导数在整段上为正则函数递增，为负则递减，恰好又得到“同增异减”。

\textbf{注意：}方法一适用得更广，因为有些严格递增的函数在个别点导数可以为 $0$，
还有些函数在某点不能求导，却仍可能单调。
\end{proofbox}
